% Problemas del capítulo 6. Generado por tools/problemas/gen.py; no editar a mano.
\section{Problemas}\label{sec:problemas-06}

Graduados: \(\star\) directo, \(\star\star\) para pensar, \(\star\star\star\) desafío. Las respuestas están en el apéndice~\ref{ap:respuestas}; conviene intentarlos antes de mirarlas.

\begin{problema}[\(\star\)]\label{pb:06-1}
Calculá con el producto cuña el área del triángulo de vértices \((0,0)\), \((6,1)\) y \((2,5)\).
\end{problema}

\begin{problema}[\(\star\)]\label{pb:06-2}
¿Están alineados los mojones \((1,2)\), \((3,5)\) y \((7,11)\)?
\end{problema}

\begin{problema}[\(\star\star\)]\label{pb:06-3}
Un lote tiene vértices \((E,N)\) en metros: \((0,0)\), \((40,0)\), \((50,30)\), \((10,40)\), en ese orden. Calculá su superficie sumando cuñas.
\end{problema}

\begin{problema}[\(\star\star\)]\label{pb:06-4}
Aplicale al lote del problema anterior la transformación \(A=\begin{pmatrix}2&1\\1&3\end{pmatrix}\). ¿Qué superficie tiene la imagen? ¿Se invierte el sentido del recorrido?
\end{problema}

\begin{problema}[\(\star\star\star\)]\label{pb:06-5}
Para \(A=\begin{pmatrix}1&2\\2&4\end{pmatrix}\), calculá \(\det A\), el núcleo y la imagen. ¿En qué se convierte el cuadrado de lados \(e_1\) y \(e_2\)?
\end{problema}
